\documentclass{jsarticle}
\usepackage[dvipdfmx, final]{graphicx}
\usepackage{amsmath}
\usepackage{amssymb}
\usepackage{chemfig}
\begin{document}
\title{無と時間、空間、そしてものごとの相対性}
\author{}
\date{}
\maketitle

  $$ {d \over df}F + \int Cdx_m = e^{-f} + e^{f} $$
  $$ =  2 (\cos (i x \log x) - i \sin(i x \log x)) $$
 $$ \int f(x)dx = \int \Gamma(\gamma)^{'}dx_m $$
   $$ = 2 (\cos (i x \log x) - i \sin(i x \log x)) $$
   $$ \left({{\int f(x) dx} \over \log x}\right) = \lim_{\theta \to \infty}
    \left({{\int f(x) dx} \over \theta}\right) = 0,1 $$
    $$ e^{i \theta} = \cos \theta + i \sin \theta $$
    $$ \left({\int f(x) dx}\right)^{'} = 2 (i \sin(i x \log x) - \cos (i x \log
    x)) $$
    $$ = 2(- \cos(i x \log x) + i \sin(i x \log x)) $$
    $$ (\cos (i x \log x) - i \sin(i x \log x))^{'} $$
    $$ = {d \over {d {e^{i\theta}}}}\left((\cos, - \sin )\cdot
    (\sin, \cos)\right)  $$
    $$ \int \Gamma(\gamma)^{''}dx_m = \left({\int \Gamma(\gamma)^{'}}dx_m
    \right)^{\nabla L} = \left({{\int \Gamma dx_m} \cdot {{d \over d\gamma}
    \Gamma}}\right)^{\nabla L} \le 
    \left({\int \Gamma dx_m + {d \over d\gamma}\Gamma}\right)^{\nabla L} \le 
    \left({e^{f} - e^{-f}} \le {e^{-f} + e^{f}}\right)^{'}  $$
    $$ = 2 $$
    萬谷先生が、ガンにかかっていながらうつになっている人を専門に
    しているのを、深谷賢治先生のつぎに、この上の式が、
    ヒッグス場の方程式が、健康体であり、ガンマ関数の大域的微分多様体が
    ガンの真逆のパールパティーの永遠の命の方程式であり、
    ヒッグス場の方程式がガンマ関数の大域的微分多様体より大であると、
    健康体であり、小であるとガンとうつが表されている生命工学の式と
    わかり、彩日記を読んでの発見・発明の式であると言えて、
    彩さん、広島大学医学部と一緒にパールパティーの母と同じく
    5千歳行きましょう。

\end{document}
